\documentclass[11pt]{article}
\usepackage[margin=1in]{geometry}
\usepackage{amsmath,amssymb,amsthm}
\usepackage{mathtools}
\usepackage{booktabs}
\usepackage{array}
\usepackage{hyperref}
\usepackage{enumitem}

\newtheorem{theorem}{Theorem}[section]
\newtheorem{lemma}[theorem]{Lemma}
\newtheorem{corollary}[theorem]{Corollary}
\newtheorem{proposition}[theorem]{Proposition}
\newtheorem{definition}[theorem]{Definition}
\newtheorem{remark}[theorem]{Remark}

\newcommand{\HH}{\mathcal{H}}
\newcommand{\OO}{\mathcal{O}}
\newcommand{\PP}{\mathcal{P}}
\DeclareMathOperator{\diag}{diag}

\title{Rank-one immunity of superconducting pairing:\\
spectral gap stability under Coulomb repulsion\\
in the Eliashberg framework}

\author{}
\date{September 16, 2026 --- internal draft, branch \texttt{p4-conj2}, UNMERGED}

\begin{document}
\maketitle

\begin{abstract}
Within the isotropic, single-band Eliashberg (Kresin--Aslamazov--...; ``KAY'')
representation of the linearized pairing eigenvalue problem, we prove that no
repulsive contribution to the pairing (``$\Delta$'') channel --- of arbitrary
strength, spectral shape, or rank --- can create a near-degeneracy at the top
of the spectrum. Writing the linearized operator as
$M=-L_1[\varphi_Z]+P[\varphi_\Delta]$ with $\varphi_Z,\varphi_\Delta$ built from
positive boson spectral measures $f_+,f_-$ in the Berk--Schrieffer decomposition,
we show
\[
 M \;\preceq\; \HH_+ + \rho(f_-)\,e_0 e_0^{\mathsf T},\qquad \HH_+:=\OO[f_+],
\]
where $\rho(f_-)\in[0,f_-(0))$ is an explicit constant depending only on the
repulsive spectral measure, and $e_0$ is the first Matsubara basis vector. This
places the leak caused by repulsion into a \emph{rank-one} perturbation, and a
one-line duality argument then gives a uniform second-eigenvalue bound
$\lambda_2(M)\le(1-c_1^2)\mu_1+c_1^2\mu_2$, where $c_1^2=\psi_1(0)^2$ is the
weight of the leading pairing eigenvector on the first Matsubara frequency ---
an ``immunity budget'' independent of the repulsion. We give three
independent quantitative lower bounds on $c_1^2(\varpi)$ (a
Davis--Kahan $\sin\Theta$ route, a Schur-test improvement, and float64
numerical evidence) certified with outward-rounded interval arithmetic, and we
carry the underlying operator identity through a first-principles Nambu
derivation for general (non-flat) density of states, isolating the exact
convolution structure of the $Z$-channel weight and the residual (currently
open) sufficient condition under which the rank-one immunity theorem extends
beyond the flat-DOS case. All numerical claims are independently
re-derived from scratch and explicitly leveled as \emph{analytic},
\emph{certificate}, or \emph{numerical-evidence}.
\end{abstract}

\tableofcontents

%======================================================================
\section{Introduction}
%======================================================================

The two-parameter Allen--Dynes-type formula $T_c=T_c(\lambda,\langle\omega\rangle,\mu^*)$
(and its many refinements) treats the entire pairing instability of an
isotropic, single-band, electron--phonon superconductor as a function of a
\emph{single scalar}: the top eigenvalue of a linearized pairing operator at
$T=T_c$. This works remarkably well empirically, but it is a priori surprising:
Coulomb repulsion $\mu^*$ acts inside the same channel (the anomalous, or
``$\Delta$'', Nambu channel) that carries the attractive electron--phonon
pairing interaction, and there is no a priori reason a repulsive contribution
of large enough magnitude, or of an adversarially chosen spectral shape,
could not push a \emph{second} pairing mode up to near-degeneracy with the
first. If that happened, ``the'' transition temperature would no longer be
governed by a single isolated eigenvalue, and the entire scalar-$T_c$
program would be structurally fragile rather than merely empirically
successful.

This question is precisely Conjecture 2 of the KAY program as pursued in this
repository (see \texttt{theory/PROGRAM.md}; the base case $\mu^*=0$ was closed
earlier, \texttt{memory/p4-theorem-status-2026-09-12.md}): does rank-one
dominance of the top pairing eigenvalue survive at $\mu^*>0$? A priori, the
answer is not obvious even in sign: repulsion lowers the top eigenvalue
$\lambda_1$ (that much is elementary, since $M\preceq\HH_+$ when repulsion is
subtracted with a negative semi-definite contribution), but lowering
$\lambda_1$ says nothing by itself about whether $\lambda_2$ might rise to
meet it.

The present paper answers this question unconditionally, in the isotropic
single-band, flat-density-of-states setting that defines the KAY
representation: Theorem~\ref{thm:main} below shows that repulsion in the
pairing channel is a rank-one perturbation \emph{in effect}, no matter how it
is spread across Matsubara frequencies or how strong it is, so that the second
eigenvalue is trapped below a value determined entirely by the attractive
sector. We further quantify the resulting ``immunity budget''
$c_1^2=\psi_1(0)^2$ --- the weight the leading attractive eigenvector places
on the first Matsubara frequency --- with three independently derived
certificate-level and numerical-evidence lower bounds, and we extend the
underlying operator identity, via a first-principles Nambu self-energy
derivation, to general (non-flat) density of states, where the analytic proof
currently closes only under a monotonicity condition on the DOS shape factor
and remains numerically (but not yet analytically) robust more generally.

\paragraph{Organization.} Section~\ref{sec:setup} recalls the KAY operator
representation and the three structural lemmas it inherits (Bochner
positivity, even-sector compression, an elementary scalar inequality).
Section~\ref{sec:main} states and proves the main result, Theorem~I
(``$\Delta$-channel immunity''), together with its two corollaries: complete
immunity to static (instantaneous) Coulomb repulsion, and the identification
of $c_1^2$ as the immunity budget. Section~\ref{sec:c1sq} gives three
independent quantitative lower bounds on $c_1^2(\varpi)$, all cross-checked
for directional soundness against a documented sign error found and
retracted during this project. Section~\ref{sec:dos} carries the derivation
to general density of states from the bare Nambu self-energy equations,
identifies the exact (and previously mis-guessed) form of the $Z$-channel
weight, and reports the current, partially open, state of the analytic
sufficient condition together with 80-configuration numerical evidence.
Section~\ref{sec:discussion} discusses the physical reading and relation to
existing KAY results. Appendix~\ref{app:iv} collects the interval-arithmetic
technicalities used to certify the numerical bounds of
Section~\ref{sec:c1sq}.

\paragraph{Rigor bookkeeping.} Every claim in this paper is tagged
\textsc{[analytic]}, \textsc{[certificate]}, or \textsc{[numerical-evidence]}
at the point it is made, following the working convention of this research
program: \textsc{analytic} means proved in closed form; \textsc{certificate}
means established by outward-rounded interval arithmetic (Sylvester inertia
counts, interval Rayleigh quotients); \textsc{numerical-evidence} means a
finite-grid float64 computation. No claim is silently promoted between these
levels.

%======================================================================
\section{Setup and notation}\label{sec:setup}
%======================================================================

We work in $\ell^2(\mathbb N_0)$, Matsubara index $n=0,1,2,\dots$, and set
\[
 u_n:=(2n+1)^{-1/2},\qquad D_u:=\diag(u_0,u_1,\dots),\qquad e_0:=(1,0,0,\dots)^{\mathsf T}.
\]
For a scalar sequence $f:\mathbb N_0\to\mathbb R$ (a boson kernel, i.e.\ a
Matsubara-frequency-difference propagator, in units where the temperature
sets $2\pi T=1$) define three operators on $\ell^2(\mathbb N_0)$:
\begin{align}
 L_1[f] &:= \diag\!\Big(\tfrac{1}{2n+1}\cdot 2\textstyle\sum_{k=1}^n f(k)\Big)
 &&\text{(mass/$Z$-channel, diagonal)},\label{eq:L1}\\
 L_2[f]_{nm} &:= (1-\delta_{nm})\,f(|n-m|)\,u_n u_m, &&\\
 L_3[f]_{nm} &:= f(n+m+1)\,u_n u_m, &&\\
 P[f] &:= L_2[f]+L_3[f] &&\text{($\Delta$/pairing channel)},\\
 \OO[f] &:= -L_1[f]+P[f], \qquad \OO'[f]:=+L_1[f]+P[f]. &&
\end{align}
($L_1$ is normalized so that $L_1[f]_{00}=0$: the $n=0$ sum is empty.)

The linearized isotropic single-band Eliashberg pairing eigenvalue problem, at
$T=T_c$, takes the KAY form
\begin{equation}\label{eq:M}
 M = -L_1[\varphi_Z] + P[\varphi_\Delta],
\end{equation}
where $\varphi_Z\ge 0$ is the total spectral weight entering the mass
renormalization ($Z$) channel and $\varphi_\Delta$ is a signed combination
entering the pairing ($\Delta$) channel. In the Berk--Schrieffer
decomposition of an attractive phonon-mediated interaction plus a repulsive
(screened Coulomb / spin-fluctuation) interaction, one writes
$\varphi_Z=f_++f_-$, $\varphi_\Delta=f_+-f_-$, with $f_+,f_-$ both drawn from
the class
\begin{equation}\label{eq:P-class}
 \PP:=\Big\{\,f(k)=\int_0^\infty \frac{\omega^2}{\omega^2+k^2}\,dP(\omega)\;:\;P\ge0\text{ a finite measure}\,\Big\}
\end{equation}
of (Einstein-mode superpositions of) bosonic exchange kernels ($\omega$
measured in units of $2\pi T$). This class is exactly the set of boson
spectral functions consistent with a non-negative Eliashberg function
$\alpha^2F\ge0$.

We record, without re-proof (this is the KAY construction; see
\texttt{theory/derivations/P4\_conjecture2.md} and the Conjecture-2 closure at
$\mu^*=0$, memory note \texttt{p4-theorem-status-2026-09-12}), that
$u_n=(2n+1)^{-1/2}$ arises from the similarity transform that symmetrizes the
linearized gap equation ($\Delta_n=\psi_n/\sqrt{2n+1}$), and is at this stage
a purely \emph{kinematic} factor (a fact re-derived independently from the
Nambu self-energy equations in Section~\ref{sec:dos}).

\subsection{Three inherited structural lemmas}

The following three lemmas are proved from scratch in this project (they do
not require the flat-DOS assumption of Eq.~\eqref{eq:M}, only that $f\in\PP$);
they are the analytic backbone of Theorem~\ref{thm:main}.

\begin{lemma}[Bochner/Poisson positivity]\label{lem:bochner}
\textup{\textsc{[analytic]}} For $f\in\PP$, extend $f$ to $\mathbb Z$ by
$f(-k)=f(k)$; the resulting sequence is positive-definite, with symbol
\begin{equation}
 \hat f(\theta):=\sum_{k\in\mathbb Z} f(k)\,e^{-ik\theta}
 =\int_0^\infty \pi\omega\,\frac{\cosh(\omega(\pi-\theta))}{\sinh(\pi\omega)}\,dP(\omega)\;\ge0,
 \qquad \theta\in[0,2\pi],
\end{equation}
using the elementary identity
$\sum_{n\in\mathbb Z}e^{-in\theta}\,\dfrac{2a}{a^2+n^2}=\dfrac{2\pi\cosh(a(\pi-\theta))}{\sinh(a\pi)}$.
Since $\cosh(\omega(\pi-\theta))$ is minimized at $\theta=\pi$,
\begin{equation}
 \min_\theta \hat f(\theta)=\hat f(\pi)=\int_0^\infty \frac{\pi\omega}{\sinh(\pi\omega)}\,dP(\omega)>0.
\end{equation}
\end{lemma}

\begin{lemma}[Even-sector compression]\label{lem:evensector}
\textup{\textsc{[analytic]}} Let $f$ be any positive-definite sequence on
$\mathbb Z$ and $x\in\ell^2(\mathbb N_0)$. Reflect $x$ about $-1/2$:
$X_n=x_n$, $X_{-1-n}=x_n$ ($n\ge0$). Then, with
$T_f=[f(n-m)]$, $H_f=[f(n+m+1)]$,
\begin{equation}
 \langle X,\tilde T X\rangle = 2\langle x,(T_f+H_f)x\rangle,\qquad \|X\|^2=2\|x\|^2,
\end{equation}
where $\tilde T$ is the (positive-definite, by hypothesis) bi-infinite
Toeplitz operator with symbol $\hat f$. Since $\langle X,\tilde TX\rangle\ge
\hat f(\pi)\|X\|^2$,
\begin{equation}
 T_f+H_f \;\succeq\; \hat f(\pi)\cdot I
\end{equation}
(and this survives restriction to any principal truncation, being a
compression statement). Consequently, since $P[f]=D_u(T_f+H_f)D_u-f(0)D_u^2$
(the $L_2$ diagonal deletion exactly removes $f(0)D_u^2$),
\begin{equation}
 P[f]\;\succeq\; -\big(f(0)-\hat f(\pi)\big)\,D_u^2 \;=:\; -\rho(f)\,D_u^2. \tag{Corollary 2.1}
\end{equation}
\end{lemma}

\begin{lemma}[Elementary scalar inequality]\label{lem:psi}
\textup{\textsc{[analytic]}} For all $\omega>0$,
\begin{equation}
 \psi(\omega):=\frac{2\omega^2}{\omega^2+1}+\frac{\pi\omega}{\sinh(\pi\omega)}-1\;\ge\;0,
 \qquad \psi(\omega)=\big(2-\tfrac{\pi^2}{6}\big)\omega^2+O(\omega^4),\quad 2-\tfrac{\pi^2}{6}=0.355066\ldots
\end{equation}
\end{lemma}
\begin{proof}
Multiplying by $(1+\omega^2)\sinh(\pi\omega)>0$, this is equivalent to
$q(\omega):=\pi\omega(1+\omega^2)-(1-\omega^2)\sinh(\pi\omega)\ge0$. Expanding
in odd powers of $\omega$, the coefficient of $\omega^{2j+1}$ is
\[
 j=0:\ \pi-\pi=0;\qquad j=1:\ 2\pi-\tfrac{\pi^3}{6}=1.115473\ldots>0;\qquad
 j\ge2:\ \frac{\pi^{2j-1}\big[2j(2j+1)-\pi^2\big]}{(2j+1)!}>0
\]
(for $j\ge2$, $2j(2j+1)\ge20>\pi^2=9.8696\ldots$). All coefficients are
non-negative, so $q(\omega)\ge0$ for all $\omega\ge0$ directly, with no
case split on $\omega$.
\end{proof}

%======================================================================
\section{Main result: Theorem I ($\Delta$-channel immunity)}\label{sec:main}
%======================================================================

\subsection{Lemma 4: the $Z$ channel dominates its own repulsive $\Delta$-channel leak}

\begin{lemma}[Rank-one leak]\label{lem:leak}
\textup{\textsc{[analytic]}} For every $f\in\PP$,
\begin{equation}\label{eq:lem4}
 \OO'[f]=L_1[f]+P[f]\;\succeq\;-\rho(f)\cdot e_0e_0^{\mathsf T},\qquad \rho(f)=f(0)-\hat f(\pi).
\end{equation}
\end{lemma}
\begin{proof}
By Corollary~2.1 (Lemma~\ref{lem:evensector}), $\OO'[f]\succeq L_1[f]-\rho(f)D_u^2=:\diag(a)$,
\[
 a_n=\frac{2\sum_{k=1}^n f(k)-f(0)+\hat f(\pi)}{2n+1}.
\]
For $n=0$ (empty sum), $a_0=\hat f(\pi)-f(0)=-\rho(f)$. For $n\ge1$, $f\ge0$
gives $\sum_{k\le n}f(k)\ge f(1)$, hence
\[
 a_n\;\ge\;\frac{2f(1)+\hat f(\pi)-f(0)}{2n+1}
 =\frac{\int\psi(\omega)\,dP(\omega)}{2n+1}\;\ge\;0
\]
by Lemma~\ref{lem:psi}. Thus $\diag(a)\succeq-\rho(f)e_0e_0^{\mathsf T}$.
\end{proof}

\begin{remark}
This is the crux of the whole argument. In an Eliashberg problem the same
piece of spectral weight $f_-$ appears simultaneously in the $Z$ channel
(positive-definite, diagonal, $L_1$) and in the $\Delta$ channel ($P$). The
diagonal growth $2\sum_{k\le n}f(k)/(2n+1)$ of $L_1$ \emph{pointwise}
dominates the most negative part $-\rho(f)u_n^2$ of $P$; the only escapee is
$n=0$, where the sum is empty. That escaped cell is rank one.
\end{remark}

\subsection{Statement and proof of Theorem I}

\begin{theorem}[$\Delta$-channel immunity]\label{thm:main}
Let $f_+,f_-\in\PP$, and let $\varphi_Z\ge0$ satisfy
$\sum_{k\le n}\varphi_Z(k)\ge\sum_{k\le n}(f_++f_-)(k)$ for all $n\ge1$
(with equality in the pure Eliashberg case $\varphi_Z=f_++f_-$). Write
$\HH_+:=\OO[f_+]$ with spectrum $\mu_1>\mu_2\ge\cdots$, top eigenvector
$\psi_1$, $c_1:=|\langle\psi_1,e_0\rangle|$, $\gamma:=\mu_1-\mu_2$. Then
\begin{equation}\label{eq:majorant}
 M\;\preceq\;\HH_++\rho(f_-)\cdot e_0e_0^{\mathsf T},\qquad \rho(f_-):=f_-(0)-\hat f_-(\pi)\in[0,f_-(0)),
\end{equation}
so that $\lambda_{j+1}(M)\le\lambda_j(\HH_+)$ for all $j\ge1$, and
\begin{equation}\label{eq:lambda2bound}
 \lambda_2(M)\;\le\;\mu_1-c_1^2\gamma \;=\; (1-c_1^2)\mu_1+c_1^2\mu_2,
\end{equation}
uniformly in $\rho$ --- i.e.\ independently of the strength, spectral shape,
or rank of the repulsion.
\end{theorem}
\begin{proof}[Proof \textup{\textsc{[analytic]}}]
The part-sum hypothesis on $\varphi_Z$ compares two diagonal operators
entrywise, giving $L_1[\varphi_Z]\succeq L_1[f_++f_-]$. Hence
\[
 M=-L_1[\varphi_Z]+P[f_+]-P[f_-]
 \;\preceq\; -L_1[f_+]-L_1[f_-]+P[f_+]-P[f_-]
 = \HH_+-\OO'[f_-]
 \;\preceq\; \HH_++\rho(f_-)\,e_0e_0^{\mathsf T},
\]
the last step by Lemma~\ref{lem:leak}. This proves Eq.~\eqref{eq:majorant}.
Weyl monotonicity gives $\lambda_j(M)\le\lambda_j(\HH_++\rho e_0e_0^{\mathsf T})$;
Cauchy interlacing for the rank-one perturbation gives
$\lambda_{j+1}(\HH_++\rho e_0e_0^{\mathsf T})\le\lambda_j(\HH_+)$.

For the uniform bound on $\lambda_2$: write $B:=\HH_++\rho e_0e_0^{\mathsf T}$,
$x:=\lambda_2(B)$, $c_j:=\langle\psi_j,e_0\rangle$. If $x\le\mu_2$ then
$x\le\mu_1-\gamma\le\mu_1-c_1^2\gamma$ (since $c_1^2\le1$) and we are done.
Otherwise $x\in(\mu_2,\mu_1)$ and the secular equation
$1=\rho\sum_j c_j^2/(x-\mu_j)$ holds. Set $d:=\mu_1-x>0$. Then
\[
 \frac{\rho c_1^2}{d}
 = \rho\sum_{j\ge2}\frac{c_j^2}{x-\mu_j}-1
 \;\le\; \frac{\rho(1-c_1^2)}{x-\mu_2}
 = \frac{\rho(1-c_1^2)}{\gamma-d}
\]
(every $j\ge2$ term has denominator $x-\mu_j\ge x-\mu_2=\gamma-d>0$, so all
terms are positive; $\sum_{j\ge2}c_j^2=1-c_1^2$; dropping the $-1$ only
enlarges the right side). Rearranging, $c_1^2(\gamma-d)\le(1-c_1^2)d$, i.e.
$d\ge c_1^2\gamma$, which is Eq.~\eqref{eq:lambda2bound}. (If $c_1=0$ then
$\lambda_2(B)=\mu_1$ and the bound holds with equality; the case that $x$
coincides exactly with some eigenvalue of $\HH_+$ is handled identically by
restricting to the corresponding invariant subspace.)

Finally, when $f_+>0$ the off-diagonal entries $f_+(|n-m|)u_nu_m>0$ of
$P[f_+]$ make $\HH_+$ irreducible, and Perron--Frobenius gives $\psi_1>0$
pointwise, hence $c_1>0$ and $\mu_1>\mu_2$ strictly.
\end{proof}

\begin{remark}[Physical reading]
Repulsion in the pairing channel can do whatever it likes there --- arbitrary
strength, arbitrary spectral shape, arbitrary rank --- and the second
eigenvalue is capped by a number determined entirely by the attractive
sector. Repulsion can suppress $\lambda_1$, but it cannot manufacture a
near-degeneracy. Rank-one (single-mode) dominance is immune to $\Delta$-channel
repulsion.
\end{remark}

\subsection{Corollaries}

\begin{corollary}[Static Coulomb: complete immunity]\label{cor:static}
\textup{\textsc{[analytic]}} Instantaneous Coulomb repulsion enters only the
$\Delta$ channel and is rank-one positive semi-definite in the KAY
representation: $P_C=2\tilde\mu\, u_c u_c^{\mathsf T}\succeq0$
($u_c[n]=u_n\mathbf 1_{n<N_c}$). It has no $Z$-channel partner, but needs
none:
\[
 M=\HH_+-2\tilde\mu\,u_cu_c^{\mathsf T}\preceq\HH_+ \quad\Longrightarrow\quad \lambda_2(M)\le\mu_2(\HH_+).
\]
Coulomb repulsion \emph{cannot} raise the second eigenvalue at all; the only
thing it can do is push $\lambda_1$ down towards $\mu_2(\HH_+)$ via Cauchy
interlacing --- killing pairing, but never manufacturing near-degeneracy.
This is the structural reason behind the numerically observed
$\mathrm{gap}/\lambda_1\to1.001$ near the $\kappa\approx2$ ``near-crossing''
point studied in the P4D Coulomb-uniqueness thread (see
\texttt{memory/negative-kernel-gap-collapse-is-z-channel-artifact-2026-09-16}).
\end{corollary}

\begin{corollary}[Immunity budget $=c_1^2$]\label{cor:budget}
\textup{\textsc{[analytic + numerical-evidence]}} Define $T_c$ by
$\lambda_1(M)=1$, and let $\mu_1=\lambda_1(\HH_+)$ at the \emph{same}
temperature (repulsion lowers $T_c$, so $\mu_1>1$). Define the suppression
share
\[
 s:=1-\frac1{\mu_1}\in[0,1) \qquad\text{(the fraction of the pairing
 eigenvalue eaten by repulsion at $T_c$)}.
\]
Then, by Theorem~\ref{thm:main} (with $\mu_2\approx0$),
\[
 \lambda_2(M)\;\le\;(1-c_1^2)\mu_1+c_1^2\mu_2\;<\;1=\lambda_1(M)
 \quad\Longleftarrow\quad s<c_1^2 \quad(\text{necessary and sufficient at }\mu_2=0),
\]
\[
 \mathrm{gap}/\lambda_1 \;\ge\; 1-(1-c_1^2)\mu_1-c_1^2\mu_2.
\]
$c_1^2$ --- the weight of the top pairing mode at the first Matsubara
frequency --- \emph{is} the immunity budget: repulsion can eat at most a
share $c_1^2$ of the pairing eigenvalue before the isolation of the top mode
could even in principle be threatened.
\end{corollary}

Table~\ref{tab:c1sq-baseline} gives the \textsc{numerical-evidence} baseline
for $c_1^2(\varpi)$ and the resulting threshold $\mu_1^{\max}=1/(1-c_1^2)$, for
the single-Einstein-mode kernel $f(k)=\varpi^2/(\varpi^2+k^2)$ (float64
\texttt{eigvalsh}, agreement to four digits between $N=400$ and $N=800$; a
$\sim10^{-2}$ truncation noise appears in $\mu_2$ for $\varpi\gtrsim30$, not
affecting $c_1^2$).

\begin{table}[h]
\centering
\caption{Baseline (numerical-evidence) values of $\mu_1$, $c_1^2$, and the
repulsion threshold $\mu_1^{\max}=1/(1-c_1^2)$, single-Einstein-mode kernel.}
\label{tab:c1sq-baseline}
\begin{tabular}{lcccccccc}
\toprule
$\varpi$ & 0.5 & 1 & 2 & 3 & 5 & 10 & 30 & 100\\
\midrule
$\mu_1$ & 0.2697 & 0.7158 & 1.3360 & 1.7177 & 2.2048 & 2.8735 & 3.9465 & 5.133\\
$c_1^2$ & 0.828 & 0.795 & 0.719 & 0.660 & 0.585 & 0.498 & 0.397 & 0.322\\
$\mu_1^{\max}$ & 5.8 & 4.9 & 3.6 & 2.9 & 2.4 & 2.0 & 1.66 & 1.48\\
\bottomrule
\end{tabular}
\end{table}

In the physically relevant range $\varpi_c\approx1$--$5$ (with $\mu^*$ or spin
fluctuations pushing $\mu_1$ to $1.2$--$1.8$), the suppression share is
$s\approx0.17$--$0.45$ against a budget $c_1^2\approx0.59$--$0.80$: the entire
physical domain falls inside the immunity budget with substantial margin.

\subsection{Escape routes: what the theorem does not cover}

Each hypothesis of Theorem~\ref{thm:main} corresponds to exactly one way the
mechanism could in principle fail (Table~\ref{tab:escape}).

\begin{table}[h]
\centering
\caption{Escape-route audit for Theorem~\ref{thm:main}.}
\label{tab:escape}
\begin{tabular}{p{4.6cm}p{5cm}p{2.6cm}}
\toprule
Violated hypothesis & Physical meaning & Physically realizable?\\
\midrule
$\sum\varphi_Z\ge\sum(f_++f_-)$ fails ($Z$-channel negative weight) & $Z<1$, negative mass renormalization & \textbf{No} ($\alpha^2F\ge0\Rightarrow$ $Z$-channel weight $\ge0$)\\
$f_-\notin\PP$ & Repulsive spectral density changes sign & \textbf{No} (spectral density $\ge0$)\\
$u_n=(2n+1)^{-1/2}$ itself changes & DOS structure: van Hove makes $\tau_n$ $n$- and $T$-dependent & \textbf{Yes} --- the only physical escape route\\
$M$ no longer of this form & Anisotropy / multi-band / momentum structure & Yes (but that is a different operator)\\
\bottomrule
\end{tabular}
\end{table}

That is: \emph{in the isotropic single-band Eliashberg setting, the only
physical channel that could break rank-one dominance is the density of
states / $Z$ channel}. This is taken up in Section~\ref{sec:dos}. An earlier
purported counterexample (a ``Mexican-hat'' kernel driving
$\mathrm{gap}/\lambda_1\to0.17$) was retracted upon audit: the entire collapse
traced to a term $+b\cdot L_1[\varphi_-]$, i.e.\ negative mass
renormalization injected directly into the $Z$ channel --- exactly the first,
non-physical, row of Table~\ref{tab:escape}
(see \texttt{theory/notes/rank\_collapse\_correction\_2026-09-16.md} and the
triage note \texttt{theory/notes/triage\_20260916\_batch\_2026-09-16.md}).
Theorem~\ref{thm:main} retroactively explains why this construction is
necessary for any such collapse to occur: no amount of negative weight placed
\emph{purely} in the $\Delta$ channel can reproduce it.

\subsection{Numerical corroboration}

\textsc{[numerical-evidence]}, from-scratch re-implementations, scripts in
\texttt{theory/notes/\_scratch\_ir/} (\texttt{probe\_rank1\_majorant.py},
\texttt{probe\_psi\_and\_c1.py}, \texttt{verify\_immunity\_chain\{,2\}.py}),
not committed, not importing repository modules:

\begin{itemize}[nosep]
\item Lemma~\ref{lem:psi}: $\min_{\omega>0}\psi(\omega)=3.55\times10^{-13}$ as
$\omega\to0$; $\psi/\omega^2\to0.355066=2-\pi^2/6$; series coefficients
positive for $j\le119$ (\texttt{mpmath} dps$=60$, series vs. closed form
relative error $\le1.5\times10^{-10}$).
\item Lemma~\ref{lem:evensector}: $\lambda_{\min}(T+H)-\hat f(\pi)$ ranges
from $+4.6\times10^{-6}$ to $-3.1\times10^{-14}$ ($\varpi=0.2,1,5,50$;
$N=200,400$) --- the bound is essentially tight, approached from above.
\item Lemma~\ref{lem:leak}: $\lambda_{\min}(\OO'[f]+f(0)e_0e_0^{\mathsf T})\ge
+2.0\times10^{-6}$ over 12 random multi-peak boson measures (1--4 peaks,
$\omega\in[0.06,50]$); over the same set $\lambda_{\min}(\OO')$ dips to
$-0.40$, confirming the rank-one correction is load-bearing.
\item Theorem~\ref{thm:main}: over 20 random $(f_+,f_-,\text{extra }Z\text{
weight})$ configurations, zero violations of $\lambda_2(M)\le(1-c_1^2)\mu_1+c_1^2\mu_2$;
the actual $\lambda_2(M)$ is almost always negative ($-0.55$ to $+0.0005$),
i.e.\ the bound is loose by roughly two orders of magnitude.
\item Rank-one secular bound (uniform in $\rho$): for $\rho=0.1$ to $10^6$,
$\varpi=1,5,30$: $\lambda_2$ saturates to $0.0134/0.495/2.013$ against bounds
$0.147/0.916/2.375$.
\end{itemize}

\emph{Directional audit.} Every numerical check above verifies whether an
\emph{upper} bound is violated; by the one-sided-bound direction rule, these
runs can only produce PASS (support), never a FAIL verdict --- a violation
would indicate a bug in the proof, which did not occur. The theorem itself
does not rely on any of these numbers.

\emph{Looseness.} The actual $\lambda_2(M)$ is essentially always negative
while the bound $(1-c_1^2)\mu_1$ is positive --- loose by roughly two orders
of magnitude. The looseness is located: (a) Lemma~\ref{lem:leak} retains only
the $n=0$ leak, while the actual negativity of $P[f_-]$ along the top mode of
$\HH_+$ is far smaller than $\rho(f_-)$; (b) the secular bound saturates only
as $\rho\to\infty$. This is an acceptable price for a ``no-go'' result: it
buys complete uniformity in $f_-$.

%======================================================================
\section{Certificate-level quantification of $c_1^2$}\label{sec:c1sq}
%======================================================================

Corollary~\ref{cor:budget} makes $c_1^2(\varpi)$ the single quantity that
converts Theorem~\ref{thm:main} into a physical statement (``the immunity
budget covers the physical domain''). Table~\ref{tab:c1sq-baseline} gives
this at \textsc{numerical-evidence} level only. This section develops three
independent lower bounds of increasing rigor, all for the single-Einstein-mode
kernel $f(k)=\varpi^2/(\varpi^2+k^2)$, $\HH_+=-L_1[f]+L_2[f]+L_3[f]$, with
top eigenvalue $\mu_1>\mu_2\ge\dots$, top eigenvector $\psi_1$, and
$c_1=|\langle\psi_1,e_0\rangle|$.

\subsection{A retracted route, kept for the record: a genuine sign error}\label{sec:c1sq-fail}

\textsc{[analytic --- retracted]}. An initial attempt used a Feshbach/Schur
complement self-energy: let $H'$ be $\HH_+$ restricted to $n\ge1$ (row/column
0 deleted), $b_n=(\HH_+)_{n0}$ ($n\ge1$). The standard identity is
\[
 c_1^2=\frac1{1+\Sigma'(\mu_1)},\qquad \Sigma(z)=b^{\mathsf T}(H'-z)^{-1}b,\qquad
 \Sigma'(z)=b^{\mathsf T}(H'-z)^{-2}b.
\]
The (incorrect) reasoning assumed Cauchy interlacing gives
$\mu_1(H')\le\mu_2(\HH_+)$, hence $\mu_1-\lambda_j(H')\ge\gamma$ for all $j$,
giving $\Sigma'(\mu_1)\le\|b\|^2/\gamma^2$ and
$c_1^2\ge\gamma^2/(\gamma^2+\|b\|^2)$. \textbf{This is backwards.} For an
$(n+1)\times(n+1)$ matrix with one row/column deleted, the correct interlacing
is $\lambda_1(H)\ge\lambda_1(H')\ge\lambda_2(H)\ge\lambda_2(H')\ge\cdots$,
i.e.\ $\mu_1(H')\ge\mu_2(\HH_+)$ (a lower bound, not an upper bound), and
$\mu_1(H')$ can sit arbitrarily close to $\mu_1$ itself --- the whole route
is circular (it would need to bound exactly the quantity, $c_1$, that it is
trying to prove a bound for). This was caught by the direction audit: the
(wrong) formula gave $(c_1^2)_{\text{lower}}=0.473$ at $\varpi=30$, $N=500$,
exceeding the numerical-evidence value $0.397$ --- a lower bound exceeding
the true value is a contradiction, flagged and discarded before any number
from this route was used further.

\subsection{Route 1: Davis--Kahan $\sin\Theta$}\label{sec:dk}

\textsc{[certificate]}, independently re-derived and re-verified
(\texttt{theory/notes/\_scratch\_ir/c1sq\_lower\_bound.py} and
\texttt{c1sq\_audit.py}), full statement in
\texttt{theory/notes/c1sq\_lower\_bound\_2026-09-16.md}.

\paragraph{Decomposition.} Fix a truncation $N$; let $A:=H_N\oplus0$ (the
$(N{+}1)\times(N{+}1)$ truncated matrix embedded in $\ell^2(\mathbb N_0)$,
remaining coordinates zero), $E:=\HH_+-A$, with block structure
$E=\begin{psmallmatrix}0&C\\ C^{\mathsf T}&H^{BB}\end{psmallmatrix}$, $C$ the
head--tail coupling block, $H^{BB}$ the tail block.

\paragraph{Analytic (non-truncated) tail bounds.} Using the same Poisson
technique as Lemma~\ref{lem:bochner} (taking the max instead of the min):
$\hat f(0)=\pi\varpi\coth(\pi\varpi)=\sum_{k\in\mathbb Z}f(k)$, and
$0\preceq T_f+H_f\preceq\hat f(0)\cdot I$ (compression preserves this bound
for any subspace). With $S:=\sum_{k\ge1}f(k)=(\hat f(0)-1)/2$,
\begin{align}
 \|H^{BB}\| &\le \frac{\hat f(0)-1}{2N+3}+u_{N+1}^2(\hat f(0)+1) &&\text{[$L_1$ tail + $(L_2{+}L_3)$ tail]},\\
 \|C\| &\le u_{N+1}\cdot\frac{\hat f(0)-\hat f(\pi)}{2} &&\text{[head--tail coupling block]},
\end{align}
using the elementary self-adjoint block-norm fact ($M$ self-adjoint with
spectrum in $[-\kappa,\kappa]$ $\Rightarrow$ $\|M_{AB}\|\le\kappa$, applied
after shifting $T_f+H_f$ by $(\hat f(0)+\hat f(\pi))/2$ to symmetrize the
spectrum). Both bounds are genuinely infinite-dimensional upper bounds, not
truncated-spectrum interval enclosures; $f,u_n,\hat f(0),\hat f(\pi)$ are all
closed form, evaluated with \texttt{mpmath.iv}, so the only error is
outward-rounded roundoff.

\paragraph{A self-contained $\sin\Theta$ inequality.} Let $\varphi$ be the
exact top eigenvector of $A$, eigenvalue $\ell_a=\lambda_1(H_N)$; let $\psi$
be the true top eigenvector of $\HH_+$, eigenvalue $\mu_1$. Write
$\psi=\cos\theta\cdot\varphi+\sin\theta\cdot\varphi_\perp$
($\varphi_\perp\perp\varphi$, unit vector). Taking the inner product with
$\varphi_\perp$:
\[
 \mu_1\sin\theta = \varphi_\perp^{\mathsf T}\HH_+\psi
 = \cos\theta\,\varphi_\perp^{\mathsf T}E\varphi + \sin\theta\,\varphi_\perp^{\mathsf T}A\varphi_\perp
 +\sin\theta\,\varphi_\perp^{\mathsf T}E\varphi_\perp
 \le \cos\theta\|E\|+\sin\theta\,\lambda_2(A)+\sin\theta\|E\|
\]
($\varphi_\perp^{\mathsf T}A\varphi=0$ since $A\varphi=\ell_a\varphi$;
$\varphi_\perp^{\mathsf T}A\varphi_\perp\le\lambda_2(A)$ by Rayleigh--Ritz;
the two $\varphi_\perp^{\mathsf T}(\cdot)$ terms are bounded by Cauchy--Schwarz).
Using $\mu_1\ge\ell_a$ (compression: $\lambda_1(H)\ge\lambda_1(H_N)$ for any
truncation, with no tail control needed),
\[
 \sin\theta\cdot(\mu_1-\lambda_2(A)-\|E\|)\le\cos\theta\|E\|\le\|E\|
 \quad\Longrightarrow\quad
 \sin\theta\le\frac{\|E\|}{\delta-\|E\|},\qquad
 \delta:=\ell_a-\lambda_2(A)=\lambda_1(H_N)-\max(\lambda_2(H_N),0).
\]
$\delta$ is the truncated matrix's \emph{own} spectral gap --- not the
direction-reversed ``gap of $H'$'' from Section~\ref{sec:c1sq-fail}. A second,
negligible ($\sim10^{-16}$) layer of $\sin\theta$ accounts for using a
floating-point approximate eigenvector as the seed for the interval
computation, itself certified via an interval residual bound. The two angles
combine by the triangle inequality, and
\[
 c_1 = |\psi[0]| \ge \cos\theta\,|v_1[0]|-\sin\theta_{\text{total}}
 \quad\Longrightarrow\quad
 c_1^2 \ge \max\big(0,\ \cos\theta\,|v_1[0]|-\sin\theta_{\text{total}}\big)^2.
\]

\paragraph{Certified $\lambda_1(H_N),\lambda_2(H_N)$.} $\lambda_1(H_N)$ lower
bound: interval Rayleigh quotient (unconditional). $\lambda_2(H_N)$ upper
bound: interval $LDL^{\mathsf T}$ (no pivoting) of $H_N-cI$, negative-pivot
count via Sylvester's inertia law; if $\ge\dim-1$ negative pivots are
certified, $\lambda_2(H_N)\le c$ (margin auto-enlarged on inconclusive
pivots). See Appendix~\ref{app:iv} for the interval-arithmetic pitfalls
avoided.

Table~\ref{tab:dk} gives the resulting certificate-level lower bounds,
$N=500$ (plus $N=1000$ at $\varpi=30$).

\begin{table}[h]
\centering
\caption{Davis--Kahan certificate lower bound on $c_1^2$, $N=500$.}
\label{tab:dk}
\begin{tabular}{lcccccccc}
\toprule
$\varpi$ & 0.5 & 1 & 2 & 3 & 5 & 10 & 30 ($N{=}500$) & 30 ($N{=}1000$)\\
\midrule
$c_1^2$ lower (certificate) & 0.6859 & 0.6577 & 0.5677 & 0.4890 & 0.3739 & 0.1961 & 0.0 & 0.0268\\
$c_1^2$ (numerical-evidence) & 0.828 & 0.795 & 0.719 & 0.660 & 0.585 & 0.498 & 0.397 & 0.397\\
$\sin\Theta_{\text{total}}$ & 0.0788 & 0.0781 & 0.0911 & 0.1083 & 0.1454 & 0.2418 & 0.7386 & 0.4106\\
$\|E\|$ & 0.0197 & 0.0518 & 0.1114 & 0.1676 & 0.2793 & 0.5586 & 1.6759 & 1.1470\\
$\delta$ & 0.2693 & 0.7146 & 1.3337 & 1.7148 & 2.2009 & 2.8687 & 3.9448 & 3.9404\\
\bottomrule
\end{tabular}
\end{table}

Every certificate lower bound is strictly below the numerical-evidence true
value (direction audit passed, in sharp contrast to
Section~\ref{sec:c1sq-fail}). At $\varpi=30$ the bound requires $N=1000$ to
become non-trivial: the coupling block $\|C\|$ decays only like
$N^{-1/2}$ (via $u_{N+1}$), while independent measurement shows the true
$\|C\|$ decays like $N^{-1}$ ($\varpi=1,30$, $N=125\to2000$, measured rate
$1.00$); this factor-of-$\sqrt N$ mismatch is the entire source of the
looseness and of the failure at $\varpi=30,N=500$ (interval $LDL^{\mathsf T}$
scales as $O(N^3)$; $N=500$ took $\sim112$s, $N=1000$ took $\sim858$s in this
project, so pushing to $N\sim4400$ needed to reach a comparable margin at
$\varpi=30$ was not attempted).

\subsection{Route 2: Schur-test improvement of $\|C\|$}\label{sec:schur}

\textsc{[certificate]}, full statement in
\texttt{theory/notes/c1sq\_schur\_improvement\_2026-09-16.md}. The
$N^{-1/2}$-vs-$N^{-1}$ mismatch of Section~\ref{sec:dk} is traced to the
global operator-norm squeeze $T_f+H_f\preceq\hat f(0)I$ discarding the
pointwise decay $f(k)\sim\varpi^2/k^2$ of the kernel. Replacing the global
bound with a Schur test recovers this information.

\paragraph{Schur test.} For any matrix $M$,
$\|M\|\le\sqrt{\sup_i R_i\cdot\sup_j K_j}$, $R_i=\sum_j|M_{ij}|$,
$K_j=\sum_i|M_{ij}|$. Applied to the coupling block $C_{nm}=u_nu_m(f(|n-m|)+f(n+m+1))$
($n\le N<m$):
\[
 R_n=u_n\sum_{m>N}u_m\big(f(m-n)+f(n+m+1)\big),\qquad
 K_m=u_m\sum_{n=0}^N u_n\big(f(m-n)+f(n+m+1)\big).
\]

\paragraph{Row sums $R_n$ (analytic tail control).} Using $u_m\le u_{N+1}$
for $m\ge N+1$,
\[
 R_n\le u_nu_{N+1}\big(\mathrm{Tail}(N{+}1{-}n)+\mathrm{Tail}(n{+}N{+}2)\big),
 \qquad \mathrm{Tail}(J):=\sum_{j\ge J}f(j).
\]
Two independent analytic bounds on $\mathrm{Tail}(J)$ are combined by taking
the minimum (each individually valid, so the minimum remains valid):
\begin{align}
\text{(a) integral comparison:}\quad
 \mathrm{Tail}(J) &\le \varpi\Big(\frac\pi2-\frac{y}{1+y^2}\Big),\quad y:=\frac{J-1}\varpi
 \quad(\text{using }\arctan y\ge y/(1+y^2)),\\
\text{(b) pointwise }1/j^2\text{ comparison:}\quad
 \mathrm{Tail}(J) &\le \varpi^2\cdot\begin{cases}2,& J\le1\\ 1/(J-1),&J\ge2.\end{cases}
\end{align}
Route (a) is purely rational plus the constant $\pi/2$ (no calls to
\texttt{iv.atan}/\texttt{iv.atan2}, a known non-enclosing routine on this
platform); route (b) avoids the transcendental constant $\pi^2/6$ altogether.
The maximum over $n=0,\dots,N$ of $R_n$ is \emph{not} attained at $n=N$ (a
float64 check shows the argmax sits at $n=0$ for $\varpi=30,100$, $N=500$),
so the certificate loops over all $n$ and takes the max of valid upper
bounds.

\paragraph{Column sums $K_m$ (exact finite sum).} For fixed $m>N$, the inner
sum over $n=0,\dots,N$ is finite and evaluated exactly by interval
arithmetic. $K_m$ is proved (not assumed) monotone decreasing in $m>N$ (each
factor $u_m,f(m-n),f(n+m+1)$ is), so $\sup_{m>N}K_m=K_{N+1}$ exactly.

\paragraph{Result.} $\|C\|_{\text{Schur}}\le\sqrt{R_{\sup}\cdot K_{N+1}}$.
Table~\ref{tab:schur} gives the resulting improved certificate.

\begin{table}[h]
\centering
\caption{Schur-test certificate lower bound on $c_1^2$, $N=500$.}
\label{tab:schur}
\begin{tabular}{lcccccccc}
\toprule
$\varpi$ & 0.5 & 1 & 2 & 3 & 5 & 10 & 30 & 100\\
\midrule
$c_1^2$ lower (Schur) & 0.8016 & 0.7763 & 0.6994 & 0.6376 & 0.5577 & 0.4587 & 0.3095 & 0.0698\\
$c_1^2$ lower (Davis--Kahan baseline) & 0.6859 & 0.6577 & 0.5677 & 0.4890 & 0.3739 & 0.1961 & 0.0 & 0.0\\
$c_1^2$ (numerical-evidence) & 0.828 & 0.795 & 0.719 & 0.660 & 0.585 & 0.498 & 0.397 & 0.322\\
$\|C\|_{\text{Schur}}$ & 0.000423 & 0.001303 & 0.002893 & 0.004486 & 0.007688 & 0.015804 & 0.076457 & 0.498660\\
$\|C\|_{\text{old}}$ & 0.0163 & 0.0455 & 0.0988 & 0.1488 & 0.2480 & 0.4960 & 1.4880 & 4.9599\\
Improvement factor & 38.5$\times$ & 34.9$\times$ & 34.2$\times$ & 33.2$\times$ & 32.3$\times$ & 31.4$\times$ & 19.5$\times$ & 9.9$\times$\\
\bottomrule
\end{tabular}
\end{table}

At $\varpi=30$ the bound is recovered at $N=500$ (from $0.0$ to $0.310$,
true value $0.397$); at $\varpi=100$ it recovers a non-trivial (if loose)
bound $0.070$ against a true value of $0.322$, where the Davis--Kahan route
gave $0$. Direction audit: all Schur bounds remain strictly below the true
value (16 configurations, $\varpi\in[0.5,100]$, $N\in\{60,200\}$, certificate
vs.\ float64 SVD truth, ratio $1.9$--$5.0\times$, all PASS).

\subsection{Route 3: an analytic lower bound}\label{sec:analytic-c1sq}

A purely analytic (interval-arithmetic-free) lower bound on $c_1^2(\varpi)$
--- e.g.\ via explicit small-matrix ($N=2$ or $3$) monotonicity, bracketing
the true value from a closed-form truncation --- was attempted but
\textbf{not achieved} within this project's time budget; see
\texttt{theory/notes/c1sq\_lower\_bound\_2026-09-16.md}, \S7 (``accepted:
not achieved''). The closest candidate considered was replacing
$\lambda_1(H_N),\lambda_2(H_N),v_1[0]$ in the Davis--Kahan chain
(Section~\ref{sec:dk}) with explicit closed-form small-$\varpi$
expansions or explicit $N=1,2$ eigenvalues, but whether this yields a
non-trivial ($>0$) bound was not verified. \emph{This route remains open};
we report it honestly here rather than substitute a fabricated result, in
keeping with the project's rigor discipline (Section 1 of the working
conventions: every claim is independently re-derived from primary sources,
and no claim is asserted without a verified source).

\subsection{Summary table and honest accounting}

\begin{table}[h]
\centering
\caption{Three-way comparison, $N=500$ (plus $N=1000$/$\varpi=30$ note).}
\label{tab:c1sq-summary}
\begin{tabular}{lccccccc}
\toprule
$\varpi$ & 0.5 & 1 & 2 & 3 & 5 & 10 & 30\\
\midrule
Davis--Kahan lower (\textsc{certificate}) & 0.686 & 0.658 & 0.568 & 0.489 & 0.374 & 0.196 & 0.0 (0.027 @ $N{=}1000$)\\
Schur lower (\textsc{certificate}) & 0.802 & 0.776 & 0.699 & 0.638 & 0.558 & 0.459 & 0.310\\
Analytic lower (\textsc{analytic}) & \multicolumn{7}{c}{not achieved --- open, Section~\ref{sec:analytic-c1sq}}\\
True value (\textsc{numerical-evidence}) & 0.828 & 0.795 & 0.719 & 0.660 & 0.585 & 0.498 & 0.397\\
\bottomrule
\end{tabular}
\end{table}

The bottleneck is entirely located: it is the estimation method for the
head--tail coupling block $\|C\|$, not the achievable truncation size $N$;
the Schur test closes most of the gap to the true value by exploiting
pointwise kernel decay instead of a global operator-norm squeeze, and the
remaining looseness (a factor of roughly $1.5$--$4.6\times$ depending on
$\varpi$) is an open, but bounded and well-characterized, target for further
tightening (e.g., a weighted Schur test using the exact geometric profile
of the top eigenvector, flagged as unexplored in
\texttt{c1sq\_lower\_bound\_2026-09-16.md}~\S6).

%======================================================================
\section{General density of states: Nambu derivation and numerical evidence}\label{sec:dos}
%======================================================================

Theorem~\ref{thm:main} is proved for the flat-DOS KAY representation, in
which $u_n=(2n+1)^{-1/2}$ is a purely kinematic factor. This section derives,
from the bare isotropic single-band Nambu--Eliashberg self-energy equations,
the exact generalization of $L_1,L_2,L_3$ to a general (non-flat) density of
states $N(\xi)$, and reports the current --- partially open --- state of the
rank-one immunity theorem in that setting.

\subsection{Bare self-energy and the shared kernel $\chi$}

Migdal's theorem, isotropy, particle--hole symmetric $N(\xi)$, the linearized
$\Delta\to0$ limit, and the on-shell approximation are inherited unchanged
from the standard Eliashberg framework (\textsc{[analytic]}, standard
results, not re-derived here; see Carbotte, \emph{Rev.\ Mod.\ Phys.}\ \textbf{62},
1027 (1990)). The Nambu self-energy is
\[
 \hat\Sigma(i\omega_n)=-T\sum_m D(\omega_n-\omega_m)\int d\xi\,N(\xi)\,\tau_3\hat G(i\omega_m,\xi)\tau_3,
 \qquad D(\omega_n-\omega_m)=\lambda(n-m),
\]
\[
 \hat G(i\omega_m,\xi)=-\frac{i\tilde\omega_m\tau_0+\xi\tau_3+\varphi_m\tau_1}{\xi^2+\tilde\omega_m^2+\varphi_m^2},
 \qquad \tilde\omega_m:=\omega_mZ_m,
\]
so that $\tau_3\hat G\tau_3=-[i\tilde\omega_m\tau_0-\xi\tau_3+\varphi_m\tau_1]/(\xi^2+\tilde E_m^2)$,
$\tilde E_m:=\sqrt{\tilde\omega_m^2+\varphi_m^2}$. Integrating over $\xi$ (the
odd term $\xi/(\xi^2+\tilde E_m^2)$ vanishes for particle--hole symmetric
$N$), and defining the shell integral
\[
 \chi(E):=\int d\xi\,\frac{N(\xi)}{\xi^2+E^2}\;\ge0,
\]
one obtains, matching $\tau_0$ and $\tau_1$ components,
\begin{equation}\label{eq:zdeq}
 \omega_n(Z_n-1)=T\sum_m\lambda(n-m)\,\tilde\omega_m\chi(\tilde E_m),
 \qquad
 \varphi_n=T\sum_m[\lambda(n-m)-\mu^*_c]\,\varphi_m\chi(\tilde E_m).
\end{equation}
\emph{This is the structural hinge of the entire section}: $\chi(\tilde E_m)$
is the \emph{same} function of the \emph{same} argument in both equations,
but it multiplies different numerators --- $\tilde\omega_m$ in the $Z$
equation, $\varphi_m$ in the $\Delta$ equation. This difference is the
source of everything that follows.

\subsection{Linearization: the two channels' effective weights diverge}

Let $\tilde E_m\to|\tilde\omega_m|$, $\chi_{\text{flat}}(E):=\pi N(0)/E$,
$g(E):=\chi(E)/\chi_{\text{flat}}(E)$ ($\to1$ in the flat-DOS limit). Note
$g(0)=1$ is a general fact (not an assumption): as $E\to0$, both $\chi(E)$
and $\chi_{\text{flat}}(E)$ are dominated by the same leading $N(0)/E$
singularity whenever $N$ is continuous and positive at the Fermi level.

\paragraph{$Z$-equation source weight.}
$\tilde\omega_m\chi(|\tilde\omega_m|)=g(\tilde\omega_m)\cdot\pi N(0)\,\mathrm{sgn}(\omega_m)$:
the $1/|\tilde\omega_m|$ singularity of $\chi_{\text{flat}}$ is \emph{exactly
cancelled} by the $\tilde\omega_m$ prefactor, and no residual
$1/(2m{+}1)$-type weight remains --- this is why the textbook flat-DOS $Z$
equation shows no explicit DOS-shape factor. For general DOS,
\[
 Z_n-1=\frac{\pi N(0)T}{\omega_n}\sum_m\lambda(n-m)\,\mathrm{sgn}(\omega_m)\,g(\tilde\omega_m).
\]

\paragraph{$\Delta$-equation source weight.}
$\chi(|\tilde\omega_m|)=g(\tilde\omega_m)\cdot\pi N(0)/|\tilde\omega_m|$: here
nothing cancels the $1/|\tilde\omega_m|$, so the flat-DOS
$1/(2m{+}1)$ weight is itself a non-trivial consequence of the DOS shell
integral. With $|\tilde\omega_m|=(2|m|{+}1)\pi T$ (leading order, $Z_m\approx1$),
\[
 Z_n\Delta_n=N(0)\sum_m[\lambda(n-m)-\mu^*_c]\,g(\tilde\omega_m)\,\frac{\Delta_m}{2m+1}
 \qquad(\text{reduces to the standard flat-DOS equation at }g\equiv1,
\]
cross-checking the entire chain of algebra against the textbook Carbotte-type
equation).

\subsection{Exact form of $P$ (the $\Delta$ channel): Ansatz A confirmed}\label{sec:ansatzA}

\textsc{[analytic]}. The $\Delta$-equation kernel
$K_{nm}=N(0)[\lambda(n-m)-\mu^*_c]g(\tilde\omega_m)/(2m{+}1)$ is asymmetric
only through the source-index weight. The similarity transform
$\Delta_m=\psi_m/\sqrt{w_m}$, $w_m:=g(\tilde\omega_m)/(2m+1)$, symmetrizes it
into $\sqrt{w_n}\lambda(n-m)\sqrt{w_m}$, exactly the $P=L_2+L_3$ structure
with
\begin{equation}\label{eq:ansatzA}
 u_n^2 := \frac{g(\tilde\omega_n)}{2n+1} \qquad (\text{``Ansatz A''}),
\end{equation}
reducing to the flat-DOS KAY substitution $u_n=(2n+1)^{-1/2}$ at $g\equiv1$.
The negative-$m$ branch maps back through $i=-1-m$ to give $\lambda(n+i+1)\sqrt{w_nw_i}$,
i.e.\ $L_3$, with the same weight. \emph{This route was originally proposed
as a plausible ansatz and is now derived from the Nambu equations.}

\subsection{Exact form of $L_1$ (the $Z$ channel): neither the original ansatz nor its first correction}\label{sec:Bn}

Two successive attempts to write down the $Z$-channel generalization of $L_1$
were made and both were falsified by a from-scratch re-derivation
(\texttt{theory/notes/nambu\_weight\_audit\_2026-09-16.md}, a self-audit of
\texttt{theory/notes/nambu\_weight\_derivation\_2026-09-16.md} and
\texttt{theory/notes/dos\_escape\_channel\_2026-09-16.md}). We report the
final, correct form and its derivation, and flag the two falsified guesses
for the historical record (they matter because a downstream numerical table
was built on the wrong one and had to be retracted).

\paragraph{Falsified guesses.} ``Ansatz A'' for $L_1$ (an overall external
factor $g_n/(2n+1)$ applied uniformly, mirroring Eq.~\eqref{eq:ansatzA}) and
``Ansatz B/C'' ($g$ attached to the boson-transfer index $|n-m|$ inside the
sum, $L_1[g,f]_n=\frac2{2n+1}\sum_{k=1}^ng(\tilde\omega_k)f(k)$) both
mis-assign the index on which $g$ acts: the source (fermionic) index $m$ in
Eq.~\eqref{eq:zdeq} was silently swapped for the boson-transfer index in the
telescoping step, and this substitution is only valid when $g\equiv1$.

\paragraph{Exact form.} Carrying out the original Matsubara sum directly:
for $m\ge0$, sign $=+1$, energy label $=m$, kernel argument $|n-m|$; for
$m\le-1$, setting $i=-1-m\ge0$, sign $=-1$, energy label $=i$, kernel
argument $n+1+i$. Hence
\begin{equation}\label{eq:Bn}
 L_1[g,f]_n = \frac{B_n}{2n+1},
 \qquad
 B_n = \underbrace{\sum_{m\ge0,\,m\neq n}\lambda(|n-m|)\,g_m}_{\text{convolution with }g\text{ on the source index}}
 \;-\;\underbrace{\sum_{i\ge0}\lambda(n+1+i)\,g_i}_{\text{negative-frequency reflection term}},
\end{equation}
where the $m=n$ term ($\lambda(0)g_n$) exactly cancels against the $\Delta$-channel
diagonal term for general $g$ (this cancellation is precisely the origin of
$L_1$ starting the sum at $k=1$ and $L_2$ carrying $(1-\delta_{nm})$ in the
KAY flat-DOS operator). $B_n$ is \emph{a convolution of $g$ along the
fermionic index with $\lambda$, minus a negative-frequency reflection term}
--- neither the overall-factor form nor the transfer-index form. At $g\equiv1$
the telescoping recovers $2\sum_{k=1}^n\lambda(k)$, the original KAY $L_1$.

A further consequence missed by both falsified guesses: $B_0=\sum_{k\ge1}\lambda(k)(g_k-g_{k-1})$
is generally \emph{non-zero} for non-constant $g$ (both falsified ansätze
give $0$ identically), which rescales the strength of the $n=0$ leak that
Lemma~\ref{lem:leak} isolates.

\begin{table}[h]
\centering
\caption{Cross-check of the exact $B_n$ (Eq.~\ref{eq:Bn}) against the flat
case and the falsified formula (``C''), \textsc{[numerical-evidence]},
$f(k)=1/(1+k^2)$, source-index truncation $M=4\times10^5$
(\texttt{nambu\_L1\_index\_audit.py}).}
\label{tab:Bn-check}
\begin{tabular}{lcccc}
\toprule
$g$ shape & $n$ & exact $B_n$ (Eq.~\ref{eq:Bn}) & falsified ``C'' formula & Ansatz A value\\
\midrule
flat & 1/3/10 & 1.000000/1.600000/1.963586 & identical & identical\\
step, pseudogap-type ($i<5$: $0.05$) & 10 & \textbf{1.912256} & 0.331821 & 1.963586\\
step, van-Hove-type ($i<5$: 1, else $0.05$) & 10 & \textbf{0.149509} & 1.729944 & 0.098179\\
van-Hove-decreasing & 0 & $-0.030283$ & $0$ (structurally missing term) & $0$\\
\bottomrule
\end{tabular}
\end{table}

The flat case reproduces the KAY telescoping identically (a self-check of
Eq.~\ref{eq:Bn}); away from flat, the falsified formula is off by up to a
factor of $5.8\times$.

\subsection{Sufficient condition and its correct direction}\label{sec:sufficient}

Writing $a_n:=(B_n-\rho(f)g_n)/(2n+1)$ (Corollary 2.1 is unaffected by $u_n$'s
form, so $P[g,f]\succeq-\rho(f)D_u^2$ continues to hold for any positive
diagonal weight $D_u$), the original (falsified) derivation's claimed
sufficient condition --- ``$g$ non-increasing (van-Hove-type) implies Lemma~4
literally survives; $g$ non-decreasing (pseudogap-type) is risky'' --- was
found, upon re-derivation with the correct $B_n$, to have the
\emph{direction reversed}:

\begin{itemize}[nosep]
\item \textbf{Pseudogap-type $g$ (non-decreasing) is essentially safe}: of 80
stress-test configurations (4 values of $\varpi$ $\times$ 20 shapes,
$N=300$), only one marginally negative value occurred ($\varpi=0.3$, width
3, depth 0.02, $-8.3\times10^{-4}$); the originally reported
$\min a_n=-5.14\times10^{-2}$ under this label does \emph{not} reproduce ---
it was an artifact of the falsified $L_1$ formula.
\item \textbf{van-Hove-type $g$ (non-increasing) is the dangerous direction}:
a sharp van-Hove step (width 3, depth 0.02) gives $-0.041$ at $\varpi=1$.
\item Random non-monotone $g$ breaks worst: $-1.37$, with negative $n$
extending to $43$ sites.
\end{itemize}

This reversal is also physically self-consistent: $g(0)=1$ always; a DOS
\emph{peak} at the Fermi level (van Hove) drives $g$ \emph{down} from 1
towards the background ratio $N_{\text{bg}}/N(0)<1$, while a Fermi-level
\emph{suppression} (pseudogap) drives $g$ \emph{up} above 1. The originally
reported direction had this backwards.

\emph{Directional audit:} $a_n<0$ only falsifies the sufficient-condition
route through Lemma~\ref{lem:leak}; by the one-sided-bound rule it does
\emph{not} by itself establish that the theorem (a statement about the true
spectral gap, not about the sign of $a_n$) fails.

\subsection{Direct spectral evidence}\label{sec:direct-spectral}

Testing the operator-level conclusion directly
(\texttt{nambu\_spec\_exact.py}, \texttt{nambu\_gap\_exact.py}, $N=600$):
flat, smooth van Hove, wide van Hove, and pseudogap shapes all satisfy the
rank-one-leak conclusion $\lambda_{\min}(\OO'-a_0e_0e_0^{\mathsf T})\ge0$ even
where $a_n<0$ occurs (the sufficient condition is lossy); a sharp van Hove
(width 3, depth $50\times$) at $\varpi=5$ gives $-0.343$, and random
``wild'' $g$ fails at all tested $\varpi$. The quantity that matters
physically, $\mathrm{gap}/\lambda_1$, is reported in
Table~\ref{tab:gap-general-dos}.

\begin{table}[h]
\centering
\caption{$\mathrm{gap}/\lambda_1$ under exact $B_n$, general DOS shapes,
\textsc{[numerical-evidence]}, $N=600$.}
\label{tab:gap-general-dos}
\begin{tabular}{lccccc}
\toprule
$\varpi$ & flat & van Hove (w3,d.02) & van Hove (smooth) & pseudogap (w5,d.05) & random (wild)\\
\midrule
0.3 & 0.9996 & 1.0000 & 0.9999 & 0.9758 & 0.9998\\
1.0 & 0.9994 & 1.0000 & 0.9999 & 0.9999 & 0.9998\\
5.0 & 0.9993 & 1.0000 & 0.9999 & 1.0030 & 0.9995\\
\bottomrule
\end{tabular}
\end{table}

Precisely where the sufficient-condition route (Lemma~\ref{lem:leak}'s
literal generalization) fails, the actual spectral isolation is essentially
perfect ($\mathrm{gap}/\lambda_1\approx1.0000$). Across every DOS shape
tested in this round, rank-one dominance / immunity is \emph{not} threatened
--- what breaks is the proof route, not the conclusion.

An earlier, independent 80-configuration scan under a still-different (but
also now known to be an approximation, cf.\ Section~\ref{sec:ansatzA}--\ref{sec:Bn})
parametrization reported $\mathrm{gap}/\lambda_1\in[0.976,1.003]$
(\texttt{theory/notes/dos\_escape\_channel\_2026-09-16.md}), consistent in
sign and magnitude with the exact-$B_n$ results above.

\subsection{What remains open}\label{sec:dos-open}

No analytic sufficient condition is currently known for the exact $B_n$
(Eq.~\ref{eq:Bn}) for general DOS shape $g$; the natural monotonicity
argument used in the flat-DOS proof (factoring out a common positive weight)
fails because $B_n$ mixes $g$ evaluated at the source index (inside a
convolution) with $g$ evaluated at $n$ (via $\rho(f)g_n$) in a way that
cannot be factored for a genuinely non-constant $g$. Four candidate routes
towards an analytic sufficient condition were explored
(\texttt{theory/notes/exact\_Bn\_analytic\_condition\_routes\_2026-09-16.md}):

\begin{enumerate}[nosep]
\item \emph{Perron--Frobenius / direct verification}: closes for flat DOS;
for general $g$, closes only under an additional (untested) ``slowly
varying $g$'' condition; open in general.
\item \emph{Gershgorin / Cassini ovals}: dead end. The pairing operator $P$
is long-range (the $L_3$ term couples all $n,m$), so the top eigenvector is
delocalized and Gershgorin disks for different $n$ always overlap
(numerically, the $n=0$ disk's upper bound exceeds the true $\lambda_1$ by a
factor $1.7$--$2.3\times$ and is not separated from the rest).
\item \emph{Small-$\varpi$ perturbative expansion}: dead end. Even at
$\varpi=0.3$, the maximal $L_1$ diagonal entry is $45\%$ of $\mu_1(P)$: $L_1$
and $P$ are comparable in strength at all $\varpi$, so there is no small
parameter to expand in.
\item \emph{$g$-perturbation theory} (bound $|\lambda_2(M[g])-\lambda_2(M[1])|\le C\|\delta g\|$
via an operator-norm bound on $\partial M/\partial g$ combined with a
Davis--Kahan step): unexplored; assessed as the most promising remaining
route, requiring an estimate of $\|\partial L_1/\partial g\|_{\text{op}}$ and
a check that the relevant spectral gap does not degenerate in the physical
$\varpi$ range.
\end{enumerate}

The present status is therefore: \textbf{analytic} for flat DOS
(Theorem~\ref{thm:main}); \textbf{analytic}, with the exact convolution
structure Eq.~\eqref{eq:Bn} derived from first principles, for general DOS,
but without a general sufficient condition for the rank-one leak; and
\textbf{numerical-evidence} (Table~\ref{tab:gap-general-dos}, 80+
configurations across multiple independent scripts) that the physical
conclusion --- no near-degeneracy --- survives every tested DOS shape,
including sharp van Hove singularities and deep pseudogaps.

%======================================================================
\section{Discussion}\label{sec:discussion}
%======================================================================

\subsection{Why this is physics and not repackaging}

\begin{enumerate}
\item \textbf{It explains why the two-parameter empirical formula works.}
Allen--Dynes-type $T_c(\lambda,\langle\omega\rangle,\mu^*)$ treats pairing
instability as a function of a single scalar (the top eigenvalue).
Theorem~\ref{thm:main} shows: no matter how strong or how spectrally shaped
the repulsion in the pairing channel is, the critical mode remains an
isolated single mode --- so ``parametrize by one scalar'' is not a fitting
accident, it is a consequence of the spectral structure.
\item \textbf{It gives a falsifiable negative prediction.} In isotropic
single-band Eliashberg theory there is \emph{no} repulsion-strength- or
spectral-shape-driven pairing-mode competition or level crossing. Any
observed competition between pairing channels must come from the DOS/$Z$
channel, anisotropy, or multiple bands --- it cannot come from ``Coulomb or
spin fluctuations being too strong.''
\item \textbf{It demotes the role of $\mu^*$.} $\mu^*$ only moves $\lambda_1$
along the Cauchy interlacing chain (Corollary~\ref{cor:static} gives its
exact floor, $\mu_2(\HH_+)$); it never touches the rank structure of the
spectrum.
\item \textbf{$c_1^2$ is a new dimensionless control number}: the weight of
the top pairing mode at the first Matsubara frequency, an explicit function
of $\varpi$ (Table~\ref{tab:c1sq-baseline}) = the immunity budget.
\end{enumerate}

\subsection{Relation to other results in this program}

\begin{itemize}
\item \emph{P4D Coulomb uniqueness} ($\kappa=2\tilde\mu c_1^2/\gamma$): the
same $c_1$, $\gamma$ appear there. Corollary~\ref{cor:static} supplies the
structural reason the $\kappa$-analysis needed: the rank-one negative
perturbation is pinned below $\mu_2$ by interlacing, so $\kappa$ only
controls \emph{where} $\lambda_1$ lands, never \emph{whether} near-degeneracy
occurs.
\item \emph{Concentration bound} (profile lower bound $v_n\ge cv_0u_n$,
\texttt{memory/p4-concentration-bound-2026-09-14}): that result needs a
\emph{lower} bound on the profile of $\psi_1$; the present result needs a
\emph{lower} bound on $c_1^2=\psi_1(0)^2$. If the concentration bound is made
analytic, Table~\ref{tab:c1sq-baseline} (and Section~\ref{sec:c1sq}'s
certificate program) could in principle be promoted from
numerical-evidence/certificate to fully analytic --- the one non-analytic
link remaining in this direction.
\item \emph{DOS universality} (\texttt{theory/notes/dimension\_universality\_2026-09-16.md}):
combined with Section~\ref{sec:dos}, the two results together assert that
the density of states is the unique physical escape hatch from KAY-type
universality statements, in two independent senses (breaking the
$T$-independence of $\tau_n$, and being the only channel not covered by
Theorem~\ref{thm:main}).
\end{itemize}

\subsection{Limitations}

\begin{enumerate}
\item The $c_1^2(\varpi)$ values in Table~\ref{tab:c1sq-baseline} are
float64 truncated-spectrum computations (\textsc{numerical-evidence}).
Theorem~\ref{thm:main} itself does not use these numbers (it holds for
arbitrary $\HH_+$); only the physical claim ``the domain lies inside the
budget'' depends on them, and Sections~\ref{sec:dk}--\ref{sec:schur} promote
substantial parts of the physical range to \textsc{certificate}.
\item The theorem covers only the isotropic single-band, flat-DOS, KAY-form
operator. The positive-definiteness argument for $P$ (the reflection
extension of Lemma~\ref{lem:evensector}) does not directly transfer to
multi-band or anisotropic settings.
\item $\rho(f)=f(0)-\hat f(\pi)$ comes from the \emph{worst-case sign}
($\theta=\pi$) in Lemma~\ref{lem:bochner}; for wide kernels ($\varpi\gg1$) it
is close to $f(0)$, while for soft kernels ($\varpi\ll1$) $\rho\to0$ and
immunity is essentially total. No finer optimization over the intermediate
regime has been attempted.
\item No interval-arithmetic (\textsc{certificate}-level) results exist for
the main theorem itself; only \textsc{analytic} (the theorem) and
\textsc{numerical-evidence} (Table~\ref{tab:c1sq-baseline} and the general-DOS
tables) apply there. \textsc{Certificate}-level results in this paper are
confined to Section~\ref{sec:c1sq}'s $c_1^2$ lower bounds.
\item The general-DOS sufficient condition (Section~\ref{sec:dos-open}) is
open; only numerical evidence, not a proof, currently supports immunity
beyond flat DOS.
\end{enumerate}

%======================================================================
\appendix
\section{Interval-arithmetic technicalities}\label{app:iv}
%======================================================================

The certificate-level results of Section~\ref{sec:c1sq} rely on
\texttt{mpmath.iv} outward-rounded interval arithmetic. Three known pitfalls
of this toolchain, tracked in this repository's working conventions, were
explicitly avoided:

\begin{enumerate}
\item \textbf{\texttt{iv.atan2} does not enclose.} It is never called in
this project; the arctangent tail bound in Section~\ref{sec:schur} uses the
rational lower bound $\arctan y\ge y/(1+y^2)$ instead of any interval
inverse-trigonometric routine.
\item \textbf{A padding margin must cover exactly the error it substitutes
for.} In the interval $LDL^{\mathsf T}$ inertia count
(Section~\ref{sec:dk}), when a pivot interval straddles zero the candidate
shift $c$ is inconclusive and the margin is enlarged and the count retried,
rather than the pivot's sign being guessed.
\item \textbf{\texttt{mp.pi} is lazy.} All uses of $\pi$ in interval context
are triggered via \texttt{+iv.pi} before use, so that outward rounding is
actually applied.
\end{enumerate}

Additionally, on the local \texttt{mpmath} 1.3.0 installation,
\texttt{iv.coth}/\texttt{iv.sinh}/\texttt{iv.cosh} were found to be broken in
interval context (an internal call to a non-existent \texttt{ctx.tanh});
these were replaced by hand with $\sinh x=(e^x-e^{-x})/2$-type expansions.

\paragraph{Inertia-count $\lambda_2(H_N)$ upper bound.} For a candidate
threshold $c$, form $H_N-cI$ in interval arithmetic and perform
unpivoted symmetric Gaussian elimination (interval $LDL^{\mathsf T}$),
counting pivots whose interval lies strictly below zero. By Sylvester's law
of inertia, if this count is $\ge\dim-1$, then $H_N$ has at most one
eigenvalue $\ge c$, certifying $\lambda_2(H_N)\le c$. If some pivot interval
straddles zero, the candidate $c$ is inconclusive and the margin around the
floating-point-guessed $c$ is enlarged and retried; correctness never relies
on the floating-point guess being accurate, only on the interval computation
succeeding once a valid $c$ is found.

\paragraph{Residual-based eigenvector certification.} A floating-point
(\texttt{numpy.eigh}) top eigenvector $v_1$ of $H_N$ is used only as a seed.
Its deviation from the true top eigenvector is bounded via the residual
$r:=\|H_Nv_1-\ell_a v_1\|$ (computed with interval arithmetic) and the
self-contained $\sin\theta$ inequality of Section~\ref{sec:dk}:
$\sin\theta_1\le r/(\mathrm{gap}_1-r)$, $\mathrm{gap}_1:=\lambda_1(H_N)-\lambda_2(H_N)$
(from the same certified inertia counter). Numerically $r\sim10^{-16}$
(machine-precision residual of a double-precision \texttt{eigh}),
$\mathrm{gap}_1\sim O(1)$, giving $\sin\theta_1\sim10^{-16}$ --- negligible,
but proved, not assumed.

\medskip
\noindent\textit{Reproducibility.} All numerical and certificate scripts
referenced above are from-scratch, independent implementations (they do not
import any repository module), located in
\texttt{theory/notes/\_scratch\_ir/} on branch \texttt{p4-conj2}; per this
project's convention, probe/exploratory scripts remain uncommitted, while
the operator identities and closed-form bounds they check are recorded in
the correspondingly dated notes cited throughout this paper.

\bigskip
\noindent\textbf{Sources.} All theorem statements, proofs, and numerical
tables in this paper are transcribed from, and independently re-verified
against: \texttt{theory/notes/delta\_channel\_immunity\_2026-09-16.md}
(Theorem I and its proof); \texttt{theory/notes/nambu\_weight\_audit\_2026-09-16.md}
and \texttt{theory/notes/nambu\_weight\_derivation\_2026-09-16.md} (general-DOS
Nambu derivation, exact $B_n$, and the audit correcting its sufficient
condition); \texttt{theory/notes/dos\_escape\_channel\_2026-09-16.md}
(80-configuration numerical evidence, superseded in its $L_1$-form claims by
the audit above but retained for its Ansatz-A/$P$-channel content and its
independent numerical scan); \texttt{theory/notes/c1sq\_lower\_bound\_2026-09-16.md}
(Davis--Kahan route, including the retracted Schur-complement sign error);
\texttt{theory/notes/c1sq\_schur\_improvement\_2026-09-16.md} (Schur-test
improvement); and \texttt{theory/notes/exact\_Bn\_analytic\_condition\_routes\_2026-09-16.md}
(status of candidate analytic sufficient conditions for general DOS). No
content was taken from any file matching the description ``$N=3$ explicit
matrix, analytic lower bound reaching 101--111\% of the true value in the
physical range $\varpi=1$--$5$'' requested for Section~\ref{sec:analytic-c1sq}:
no such file exists in this repository as of 2026-09-16 (verified by
directory listing and full-text search of \texttt{theory/notes/}); this is
flagged rather than fabricated, and Section~\ref{sec:analytic-c1sq} is
written to reflect the true, open status of that route.

\end{document}
